% GENERATED FILE. Edit canonical lessons, then run npm run generate:books.
% canonical-source-hash: fnv1a64:4e3326a9f19295c1
% canonical-lessons: MR-C219-dar, MR-C219-utpanna, MR-C219-vyapar, MR-C219-dukandar, MR-C219-grahak

\chapter{Rate, Income, Trade, Shopkeeper, Customer}
\label{ch:mr-rate-income-trade-shopkeeper-customer}
\hlchaptermodality{219}

\begin{chapteropening}
\textbf{By the end of this chapter:} \emph{I can say a rate, income, trade, a shopkeeper and a customer.}

Complete the last lesson of chapter 219: I can say a rate, income, trade, a shopkeeper and a customer.
\end{chapteropening}

\section[\texorpdfstring{\mr{दर} (dar)}{दर (dar)}]{\mr{दर} (dar) — a rate}

\label{lesson:MR-C219-dar}



\begin{quote}
Before the new one: say the Marathi for fog, then the Marathi for dew.
\end{quote}

\subsection*{You'll want to know: \mr{दर}}
\textbf{\mr{दर}} — \emph{dar} — \textquotedblleft{}a rate\textquotedblright{}.

\begin{cousinweb}
The origin of \textbf{\mr{दर}} (\emph{dar}), \textquotedblleft{}rate\textquotedblright{}, is not settled. CDIAL has no entry for it in this sense.
\end{cousinweb}

\begin{grammarlens}[title={What you've built}]
One more word for this chapter.
\end{grammarlens}

\subsection*{Your turn}
\begin{itemize}
\raggedright
  \item \emph{Say it:} \emph{dar}
  \item \emph{Say it:} \emph{dar}, once more
  \item \emph{Say it:} say \emph{dar}
  \item \emph{Recall:} say the Marathi for fog, then the Marathi for dew, then say \emph{dar} again
\end{itemize}

\begin{tcolorbox}[breakable,colback=teal!4,colframe=teal!35!black,title={Before you move on}]
What does \mr{दर} mean? (\textquotedblleft{}A rate\textquotedblright{}.) Say it once more.
\end{tcolorbox}
\section[\texorpdfstring{\mr{उत्पन्न} (utpanna)}{उत्पन्न (utpanna)}]{\mr{उत्पन्न} (utpanna) — income}

\label{lesson:MR-C219-utpanna}



\begin{quote}
Before the new one: say the Marathi for dew, then the Marathi for a rate.
\end{quote}

\subsection*{You'll want to know: \mr{उत्पन्न}}
\textbf{\mr{उत्पन्न}} — \emph{utpanna} — \textquotedblleft{}income\textquotedblright{}.

\begin{cousinweb}
\textbf{\mr{उत्पन्न}} (\emph{utpanna}) is Sanskrit \emph{utpanna}, \textquotedblleft{}arisen, produced\textquotedblright{}, a tatsama: \emph{ud}, \textquotedblleft{}up\textquotedblright{}, and \emph{pad}, \textquotedblleft{}to go\textquotedblright{}. CDIAL has no entry for it.
\end{cousinweb}

\begin{grammarlens}[title={What you've built}]
One more word for this chapter.
\end{grammarlens}

\subsection*{Your turn}
\begin{itemize}
\raggedright
  \item \emph{Say it:} \emph{utpanna}
  \item \emph{Say it:} \emph{utpanna}, once more
  \item \emph{Say it:} say \emph{utpanna}
  \item \emph{Recall:} say the Marathi for dew, then the Marathi for a rate, then say \emph{utpanna} again
\end{itemize}

\begin{tcolorbox}[breakable,colback=teal!4,colframe=teal!35!black,title={Before you move on}]
What does \mr{उत्पन्न} mean? (\textquotedblleft{}Income\textquotedblright{}.) Say it once more.
\end{tcolorbox}
\section[\texorpdfstring{\mr{व्यापार} (vyāpār)}{व्यापार (vyāpār)}]{\mr{व्यापार} (vyāpār) — trade}

\label{lesson:MR-C219-vyapar}



\begin{quote}
Before the new one: say the Marathi for a rate, then the Marathi for income.
\end{quote}

\subsection*{You'll want to know: \mr{व्यापार}}
\textbf{\mr{व्यापार}} — \emph{vyāpār} — \textquotedblleft{}trade\textquotedblright{}.

\begin{cousinweb}
\textbf{\mr{व्यापार}} (\emph{vyāpār}) is Sanskrit \emph{vyāpāra}, \textquotedblleft{}occupation, business\textquotedblright{}, taken over whole as a learned loan, a tatsama (CDIAL 12205). CDIAL lists the inherited Marathi form as \emph{vāvar}, \textquotedblleft{}business, activity\textquotedblright{}.
\end{cousinweb}

\begin{grammarlens}[title={What you've built}]
One more word for this chapter.
\end{grammarlens}

\subsection*{Your turn}
\begin{itemize}
\raggedright
  \item \emph{Say it:} \emph{vyāpār}
  \item \emph{Say it:} \emph{vyāpār}, once more
  \item \emph{Say it:} say \emph{vyāpār}
  \item \emph{Recall:} say the Marathi for a rate, then the Marathi for income, then say \emph{vyāpār} again
\end{itemize}

\begin{tcolorbox}[breakable,colback=teal!4,colframe=teal!35!black,title={Before you move on}]
What does \mr{व्यापार} mean? (\textquotedblleft{}Trade\textquotedblright{}.) Say it once more.
\end{tcolorbox}
\section[\texorpdfstring{\mr{दुकानदार} (dukāndār)}{दुकानदार (dukāndār)}]{\mr{दुकानदार} (dukāndār) — a shopkeeper}

\label{lesson:MR-C219-dukandar}



\begin{quote}
Before the new one: say the Marathi for income, then the Marathi for trade.
\end{quote}

\subsection*{You'll want to know: \mr{दुकानदार}}
\textbf{\mr{दुकानदार}} — \emph{dukāndār} — \textquotedblleft{}a shopkeeper\textquotedblright{}.

\begin{cousinweb}
\textbf{\mr{दुकानदार}} (\emph{dukāndār}) is \emph{dukān}, \textquotedblleft{}shop\textquotedblright{}, from Arabic \emph{dukkān}, \textquotedblleft{}shop\textquotedblright{}, through Persian (Lane, root \emph{d-k-n}), with Persian \emph{-dār}, \textquotedblleft{}holder\textquotedblright{}.
\end{cousinweb}

\begin{grammarlens}[title={What you've built}]
One more word for this chapter.
\end{grammarlens}

\subsection*{Your turn}
\begin{itemize}
\raggedright
  \item \emph{Say it:} \emph{dukāndār}
  \item \emph{Say it:} \emph{dukāndār}, once more
  \item \emph{Say it:} say \emph{dukāndār}
  \item \emph{Recall:} say the Marathi for income, then the Marathi for trade, then say \emph{dukāndār} again
\end{itemize}

\begin{tcolorbox}[breakable,colback=teal!4,colframe=teal!35!black,title={Before you move on}]
What does \mr{दुकानदार} mean? (\textquotedblleft{}A shopkeeper\textquotedblright{}.) Say it once more.
\end{tcolorbox}
\section[\texorpdfstring{\mr{ग्राहक} (grāhak)}{ग्राहक (grāhak)}]{\mr{ग्राहक} (grāhak) — a customer}

\label{lesson:MR-C219-grahak}



\begin{quote}
Before the new one: say the Marathi for trade, then the Marathi for a shopkeeper.
\end{quote}

\subsection*{You'll want to know: \mr{ग्राहक}}
\textbf{\mr{ग्राहक}} — \emph{grāhak} — \textquotedblleft{}a customer\textquotedblright{}.

\begin{cousinweb}
\textbf{\mr{ग्राहक}} (\emph{grāhak}) is Sanskrit \emph{grāhaka}, \textquotedblleft{}one who takes, buyer\textquotedblright{}, a tatsama built on \emph{grah}, \textquotedblleft{}to take\textquotedblright{}. CDIAL has no entry for it.
\end{cousinweb}

\begin{grammarlens}[title={What you've built}]
One more word for this chapter.
\end{grammarlens}

\subsection*{Your turn}
\begin{itemize}
\raggedright
  \item \emph{Say it:} \emph{grāhak}
  \item \emph{Say it:} \emph{grāhak}, once more
  \item \emph{Say it:} say \emph{grāhak}
  \item \emph{Recall:} say the Marathi for trade, then the Marathi for a shopkeeper, then say \emph{grāhak} again
\end{itemize}

\begin{tcolorbox}[breakable,colback=teal!4,colframe=teal!35!black,title={Before you move on}]
What does \mr{ग्राहक} mean? (\textquotedblleft{}A customer\textquotedblright{}.) Say it once more.
\end{tcolorbox}
